\documentclass[11pt,a4paper]{article}

% Kompiliuoti su PDFLaTeX.
\usepackage[T1]{fontenc}
\usepackage[utf8]{inputenc}
\usepackage[provide=*,lithuanian]{babel}
\usepackage{lmodern}
\usepackage{amsmath,amssymb}
\usepackage[margin=1.9cm,headheight=15pt,headsep=12pt]{geometry}
\usepackage{enumitem,multicol,array,tabularx,longtable}
\usepackage{fancyhdr}
\usepackage{microtype}
\usepackage{xcolor}
\usepackage{titlesec}
\usepackage{needspace}
\usepackage{booktabs}

\definecolor{mainblue}{HTML}{20385C}
\definecolor{muted}{HTML}{5C6470}
\definecolor{lightline}{HTML}{DCE2E8}

\setlength{\parindent}{0pt}
\setlength{\parskip}{4pt}
\setlength{\columnsep}{20pt}
\setlist[enumerate,1]{label=\textbf{\alph*)},leftmargin=1.0cm,itemsep=.4em,topsep=.35em}
\setlist[itemize]{leftmargin=1.0cm,itemsep=.18em,topsep=.2em}
\renewcommand{\arraystretch}{1.16}

\pagestyle{fancy}
\fancyhf{}
\fancyhead[L]{\small\color{muted}Kvadratinis daugianaris}
\fancyhead[R]{\small\color{muted}9 klasė}
\fancyfoot[C]{\small\color{muted}\thepage}
\renewcommand{\headrulewidth}{0.3pt}
\renewcommand{\footrulewidth}{0pt}

\titleformat{\section}{\large\bfseries\color{mainblue}}{\thesection.}{.6em}{}
\titlespacing{\section}{0pt}{1.35em}{.5em}
\titleformat{\subsection}{\normalsize\bfseries\color{mainblue}}{}{0pt}{}
\titlespacing{\subsection}{0pt}{1em}{.35em}

\newcommand{\R}{\mathbb{R}}
\newcommand{\kom}[1]{\par\smallskip{\small\color{muted}\textit{#1}}\par}
\newcommand{\uzd}[2]{\needspace{7\baselineskip}\section{#1}\textit{#2}\par\smallskip}

\begin{document}

\begin{center}
{\LARGE\bfseries\color{mainblue} Kvadratinis daugianaris}\\[4pt]
{\large Pavidalai, savybės ir vidiniai ryšiai}\\[4pt]
{\small Kompleksinis užduočių rinkinys \(\cdot\) 9 klasė}
\end{center}
\vspace{3pt}
\noindent\textbf{Mokinio kodas:}\ \rule{4cm}{0.4pt}
\hfill\textbf{Klasė:}\ \rule{2.1cm}{0.4pt}
\vspace{4pt}
\hrule height .45pt
\vspace{4pt}

\textbf{Naudinga prisiminti.} Jei $P(x)=ax^2+bx+c$ ir $a\neq 0$, tai
\[
D=b^2-4ac,\qquad h=-\frac{b}{2a},\qquad
P(x)=a(x-h)^2+m,\qquad m=-\frac{D}{4a}.
\]
Skaičius $h$ yra ekstremumo taškas, o $m$ -- ekstremumas (mažiausioji arba
 didžiausioji reikšmė, žiūrint į $a$ ženklą).

\uzd{Trys pavidalai -- vienas daugianaris}{Iš vienos išraiškos išskaitykite kuo daugiau informacijos.}
Duotas kvadratinis daugianaris
\[
P(x)=-2x^2+12x-10.
\]
\begin{enumerate}
\item Užrašykite šį daugianarį dar dviem pavidalais:
\[
P(x)=a(x-x_1)(x-x_2),\qquad P(x)=a(x-h)^2+m.
\]
\item Remdamiesi gautais pavidalais nustatykite daugianario šaknis, ekstremumo
 tašką, ekstremumą ir jo tipą (minimumas ar maksimumas).
\item Raskite didžiausiąją ir mažiausiąją daugianario reikšmę \emph{kiekviename}
 iš intervalų
\[
[2;4]\qquad\text{ir}\qquad[4;6].
\]
Kiekvienu atveju nurodykite, su kuriomis $x$ reikšmėmis šios reikšmės
 pasiekiamos.
\item Nustatykite, su kuriomis realiosiomis $x$ reikšmėmis
 $P(x)>0$, o su kuriomis $P(x)<0$.
\item Paaiškinkite, kuris iš trijų pavidalų patogiausias šaknims,
 kuris ekstremumui, o kuris koeficientams nustatyti.
\end{enumerate}

\uzd{Įrodykite, kad ženklas nesikeičia}{Įrodymas turi galioti visoms realiosioms $x$ reikšmėms.}
Duoti kvadratiniai daugianariai
\[
P(x)=2x^2-8x+11,\qquad Q(x)=-3x^2-12x-16.
\]
\begin{enumerate}
\item Įrodykite, kad $P(x)>0$ su kiekviena $x\in\R$.
\item Įrodykite, kad $Q(x)<0$ su kiekviena $x\in\R$.
\item Ar šie daugianariai turi realiųjų šaknų? Pagrįskite.
\item Ar a) ir b) įrodymams būtina apskaičiuoti diskriminantą?
 Pateikite įrodymą išskirdami pilnąjį kvadratą.
\end{enumerate}

\uzd{Atpažinkite daugianario ženklus}{Svarbu suprasti daugianario pavidalą, o ne atlikti nereikalingus veiksmus.}
Duota dešimt kvadratinių daugianarių:
\begin{multicols}{2}
\begin{enumerate}[label=\textbf{\arabic*.},leftmargin=0.8cm,itemsep=.65em]
\item $2(x-1)^2+3$
\item $(x+4)(x-2)$
\item $-2(x+3)^2$
\item $-3(x+2)^2-4$
\item $2(x-5)^2$
\item $-2(x-1)(x-6)$
\item $(x+1)^2$
\item $3(x+3)(x-1)$
\item $-4(x-2)^2$
\item $-(x+2)(x-5)$
\end{enumerate}
\end{multicols}
\begin{enumerate}
\item \textbf{Neatskliausdami ir nespręsdami kvadratinių lygčių}, kiekvienam
 daugianariui nustatykite:
\begin{itemize}
\item ar jis turi realiųjų šaknų ir, jei turi, kokios jos;
\item su kuriomis $x$ reikšmėmis daugianaris teigiamas;
\item su kuriomis $x$ reikšmėmis daugianaris neigiamas;
\item su kuriomis $x$ reikšmėmis jis lygus nuliui.
\end{itemize}
\item Nenustatydami tikslios diskriminanto reikšmės nurodykite, kurių
 daugianarių diskriminantai teigiami, neigiami arba lygūs nuliui.
 Paaiškinkite, kaip tai nustatėte.
\item Antrąjį ir šeštąjį daugianarius pažymėkime atitinkamai $P(x)$ ir $Q(x)$.
 Raskite visas $x$ reikšmes, su kuriomis:
\begin{itemize}
\item abu daugianariai teigiami;
\item pirmasis neigiamas, o antrasis teigiamas;
\item pirmasis teigiamas, o antrasis neigiamas.
\end{itemize}
\end{enumerate}

\uzd{Sukonstruokite daugianarį}{Kiekvienu atveju paaiškinkite, kaip nustatėte koeficientus.}
Kiekvienoje dalyje raskite kvadratinį daugianarį
$P(x)=ax^2+bx+c$, kur $a\neq0$, tenkinantį nurodytas sąlygas.
\begin{enumerate}
\item \textbf{Simetrija ir ekstremumas.}
 Žinoma, kad $P(1)=P(5)=9$, o mažiausioji reikšmė yra $1$.
 Sukonstruokite daugianarį.
\item \textbf{Šaknys ir viena reikšmė.}
 Daugianario šaknys yra $-2$ ir $4$. Be to, $P(0)=-16$.
 Sukonstruokite daugianarį.
\item \textbf{Diskriminantas ir ekstremumo taškas.}
 Žinoma, kad $a=1$, $D=36$, o ekstremumo taškas yra $x=2$.
 Sukonstruokite daugianarį.
\item \textbf{Daugianario ženklai.}
 Sukonstruokite daugianarį, kuris teigiamas tik tada, kai $2<x<6$,
 neigiamas, kai $x<2$ arba $x>6$, ir tenkina $P(4)=8$.
\item \textbf{Simetrija ir diskriminantas.}
 Žinoma, kad $P(1)=P(5)$, $a=2$ ir $D=32$.
 Sukonstruokite daugianarį.
\item \textbf{Šaknys ir diskriminantas.}
 Daugianario šaknys yra $-1$ ir $5$, jo diskriminantas lygus $144$,
 o aukščiausiojo laipsnio koeficientas neigiamas.
 Sukonstruokite daugianarį.
\item \textbf{Ar daugianaris nustatomas vienareikšmiškai?}
 Sukonstruokite tris skirtingus kvadratinius daugianarius, kuriems
\[
P(x)>0\quad\text{kai }x<1\text{ arba }x>5,\qquad
P(x)<0\quad\text{kai }1<x<5.
\]
Ar sąlygos vienareikšmiškai nustato daugianarį? Jei ne, pateikite
 \emph{visų} tokių daugianarių formulę.
\end{enumerate}

\uzd{Įrodykite, kad taip būti negali}{Suraskite prieštarą duotoms daugianario savybėms.}
Kiekvienoje dalyje pateikta informacija apie kvadratinį daugianarį
$P(x)=ax^2+bx+c$ ir jai prieštaraujantis teiginys. Įrodykite, kad
 teiginys negali būti teisingas. Vien kelių konkrečių reikšmių patikrinimo nepakanka.
\begin{enumerate}
\item Žinoma, kad $a>0$ ir $D<0$.
 Mokinys teigia, jog egzistuoja $x\in\R$, kuriam $P(x)<0$.
 Įrodykite, kad jis klysta.
\item Žinoma, kad $a<0$ ir $D>0$.
 Mokinys teigia, jog didžiausioji daugianario reikšmė neigiama.
 Įrodykite, kad taip būti negali.
\item Žinoma, kad $a>0$ ir $b>0$.
 Mokinys teigia, kad ekstremumo taškas yra teigiamas skaičius.
 Įrodykite, kad taip būti negali.
\item Žinoma, kad $a>0$ ir $c<0$.
 Mokinys teigia, jog diskriminantas neigiamas.
 Įrodykite, kad taip būti negali. Pateikite \emph{du skirtingus argumentus}.
\item Žinoma, kad $a>0$ ir $P(1)=P(7)$.
 Mokinys teigia, jog $P(4)>P(3)$.
 Įrodykite, kad teiginys klaidingas.
\item Žinoma, kad $a=2$ ir $D=32$.
 Mokinys teigia, kad mažiausioji daugianario reikšmė yra $3$.
 Įrodykite, kad taip būti negali.
\item Žinoma, kad $a>0$, $b>0$ ir $D>0$.
 Mokinys teigia, kad abi daugianario šaknys teigiamos.
 Įrodykite, kad taip būti negali.
\item Žinoma, kad $a<0$ ir $P(2)=P(8)$.
 Mokinys teigia, jog $P(5)<P(7)$.
 Įrodykite, kad taip būti negali.
\end{enumerate}

\clearpage
\begin{center}
{\LARGE\bfseries\color{mainblue} Atsakymai ir sprendimų gairės}\\[3pt]
{\small Ši dalis skirta savarankiškam pasitikrinimui ir mokytojui.}
\end{center}
\vspace{8pt}

\section*{1 užduotis}
\[
P(x)=-2(x-1)(x-5)=-2(x-3)^2+8.
\]
Šaknys: $x_1=1$, $x_2=5$. Ekstremumo taškas: $x=3$.
 Maksimumas: $8$ (pasiekiamas, kai $x=3$).
\begin{center}
\begin{tabular}{@{}lll@{}}
\toprule
Intervalas & Mažiausioji reikšmė & Didžiausioji reikšmė \\
\midrule
$[2;4]$ & $6$, kai $x=2$ arba $x=4$ & $8$, kai $x=3$ \\
$[4;6]$ & $-10$, kai $x=6$ & $6$, kai $x=4$ \\
\bottomrule
\end{tabular}
\end{center}
$P(x)>0$, kai $1<x<5$; $P(x)<0$, kai $x<1$ arba $x>5$.
 Išskaidytasis pavidalas patogiausias šaknims, pilnojo kvadrato pavidalas --
 ekstremumui, o išskleistasis -- koeficientams nuskaityti.

\section*{2 užduotis}
\[
P(x)=2(x-2)^2+3>0,\qquad Q(x)=-3(x+2)^2-4<0
\quad\text{su visais }x\in\R.
\]
Kadangi kvadratas neneigiamas, pirmasis daugianaris visada ne mažesnis
 už $3$, o antrasis visada ne didesnis už $-4$.
 Todėl nė vienas neturi realiųjų šaknų.

\section*{3 užduotis}
Žemiau kiekvienam daugianariui pateikta visa ženklų analizė.
 Pažymėjimas $\R\setminus\{u\}$ reiškia visus realiuosius skaičius,
 išskyrus $u$.
\begin{enumerate}[label=\textbf{\arabic*.},leftmargin=0.85cm,itemsep=.6em]
\item $D<0$, šaknų nėra. $P>0$ su visais $x\in\R$; $P<0$ niekada;
 $P=0$ niekada.
\item $D>0$, šaknys $-4$ ir $2$.
 $P>0$: $x<-4$ arba $x>2$;
 $P<0$: $-4<x<2$;
 $P=0$: $x=-4$ arba $x=2$.
\item $D=0$, šaknis $-3$ (dviguba).
 $P>0$ niekada; $P<0$: $x\neq -3$; $P=0$: $x=-3$.
\item $D<0$, šaknų nėra. $P>0$ niekada;
 $P<0$ su visais $x\in\R$; $P=0$ niekada.
\item $D=0$, šaknis $5$ (dviguba).
 $P>0$: $x\neq5$; $P<0$ niekada; $P=0$: $x=5$.
\item $D>0$, šaknys $1$ ir $6$.
 $P>0$: $1<x<6$;
 $P<0$: $x<1$ arba $x>6$;
 $P=0$: $x=1$ arba $x=6$.
\item $D=0$, šaknis $-1$ (dviguba).
 $P>0$: $x\neq -1$; $P<0$ niekada; $P=0$: $x=-1$.
\item $D>0$, šaknys $-3$ ir $1$.
 $P>0$: $x<-3$ arba $x>1$;
 $P<0$: $-3<x<1$;
 $P=0$: $x=-3$ arba $x=1$.
\item $D=0$, šaknis $2$ (dviguba).
 $P>0$ niekada; $P<0$: $x\neq2$; $P=0$: $x=2$.
\item $D>0$, šaknys $-2$ ir $5$.
 $P>0$: $-2<x<5$;
 $P<0$: $x<-2$ arba $x>5$;
 $P=0$: $x=-2$ arba $x=5$.
\end{enumerate}
\textbf{c) dalis.} Kai $P(x)=(x+4)(x-2)$, o
$Q(x)=-2(x-1)(x-6)$:
\begin{itemize}
\item abu teigiami: $2<x<6$;
\item $P(x)<0$, $Q(x)>0$: $1<x<2$;
\item $P(x)>0$, $Q(x)<0$: $x<-4$ arba $x>6$.
\end{itemize}

\section*{4 užduotis}
\begin{enumerate}
\item $P(x)=2(x-3)^2+1$.
 Vienodų reikšmių $P(1)=P(5)$ argumentai simetriški taško $3$ atžvilgiu,
 todėl $h=3$. Iš $P(1)=9$ gauname $4a+1=9$, taigi $a=2$.
\item $P(x)=2(x+2)(x-4)$.
 Iš šaknų gauname $a(x+2)(x-4)$; sąlyga $P(0)=-16$ duoda $-8a=-16$.
\item $P(x)=x^2-4x-5$.
 Kadangi $-b/(2a)=2$ ir $a=1$, turime $b=-4$;
 iš $D=b^2-4ac=36$ gauname $c=-5$.
\item $P(x)=-2(x-2)(x-6)$.
 Pagal ženklus šaknys yra $2$ ir $6$, o $a<0$.
 Iš $P(4)=8$ gauname $-4a=8$, todėl $a=-2$.
\item $P(x)=2(x-3)^2-4$.
 Iš $P(1)=P(5)$ gauname $h=3$;
 iš $m=-D/(4a)$ gauname $m=-4$.
\item $P(x)=-2(x+1)(x-5)$.
 Čia $D=a^2(x_2-x_1)^2=36a^2=144$.
 Todėl $|a|=2$, o kadangi $a<0$, tai $a=-2$.
\item Pavyzdžiui, $(x-1)(x-5)$, $2(x-1)(x-5)$ ir
 $3(x-1)(x-5)$. Visi atsakymai:
 \[
 P(x)=a(x-1)(x-5),\qquad a>0.
 \]
\end{enumerate}

\section*{5 užduotis}
\begin{enumerate}
\item Iš formulės $m=-D/(4a)$ gauname $m>0$, nes $D<0$ ir $a>0$.
 Tai minimumas, todėl $P(x)\ge m>0$ visiems $x\in\R$.
\item Kadangi $a<0$, skaičius $m=-D/(4a)>0$ yra maksimumas.
 Jis negali būti neigiamas.
\item Ekstremumo taškas $h=-b/(2a)<0$, nes $a,b>0$.
\item Pirmas argumentas: $D=b^2-4ac>0$, nes $b^2\ge0$ ir $-4ac>0$.
 Antras argumentas: $P(0)=c<0$. Jei $a>0$ ir $D<0$,
 tai pagal pilnojo kvadrato pavidalą $P(x)>0$ visiems $x$, prieštara.
\item Iš $P(1)=P(7)$ gauname $h=4$.
 Kadangi $a>0$, taške $4$ yra minimumas, todėl $P(4)<P(3)$.
\item Mažiausioji reikšmė $m=-D/(4a)=-32/8=-4$, o ne $3$.
\item Pagal Vijetos teoremą šaknų suma yra $-b/a<0$.
 Dvi teigiamos šaknys turėtų teigiamą sumą -- prieštara.
\item Iš $P(2)=P(8)$ gauname $h=5$.
 Kadangi $a<0$, taške $5$ yra maksimumas, todėl $P(5)>P(7)$.
\end{enumerate}

\end{document}
